\documentclass[11pt]{article}
\usepackage{amsmath,amssymb,amsthm,hyperref}
\newtheorem{theorem}{Theorem}
\newtheorem{lemma}{Lemma}
\title{A spectral route and modular normal form for Green's Problem 15}
\author{Research notes; no novelty claim}
\date{6 November 2026}
\begin{document}
\maketitle

\noindent\textbf{Target retired after live status check.} On 6 November 2026,
the live Conjectures.io page was checked and reports Green 15 approved on
11 September 2026, submitted by JenW1N, with increments between 1 and 125.
The attached proof uses a six-letter substitution and endpoint tags.
This is an external result, not a result of these notes. The research below
was conducted under a stale open-status assumption and does not fulfill
the user's request to solve a currently unsolved problem. Source:
\url{https://conjectures.io/results/66df312f-30b4-4331-b1fa-04be52f7f8ca/solution}.

\noindent\textbf{Status.} Green's Problem 15 has not been solved here.
The sufficient condition below has been reviewed, but no substitution satisfying
all its hypotheses has been found. The computations exclude only specified finite
families. A long square-free prefix would not establish an infinite construction.
No Lean proof or literature priority claim is made.

An additive square is a pair of adjacent nonempty blocks with the same length
and the same sum. An abelian square has adjacent blocks with identical letter
counts. For a word $u$ on $s$ letters, write $\Psi(u)\in\mathbb Z^s$ for its
column vector of letter counts and $\ell(x)=\sum_i x_i$.

\begin{theorem}[Sufficient spectral condition]
Let $h$ be a primitive non-erasing substitution with Perron eigenvalue greater
than one. Suppose a power $g=h^t$ is prolongable on $a$, and let
$W=g^\omega(a)$. Let $M$ be the column incidence matrix of $h$.
Suppose $M$ has no eigenvalues of modulus one and that the real generalized
eigenspace $E$ corresponding to eigenvalues of modulus greater than one has
dimension two. If $W$ is abelian-square-free, then its letters admit distinct
positive integer weights for which the resulting integer word is additive-square-free.
Such a word gives a positive answer to Green's Problem 15.
\end{theorem}

\begin{proof}
Replace $g$ by a further power so every image has length at least two;
primitivity and growth permit this. Write $A$ for its incidence matrix, a
power of $M$. The expanding and contracting generalized eigenspaces of $A$
are the same $E$ and $F$ as those of $M$. Let $\pi$ be projection onto $F$
along $E$, so $\mathbb R^s=E\oplus F$. The Perron eigenvector is positive
and belongs to $E$, so $\ell$ is nonzero on $E$.
In particular, $E\cap\ker\ell$ is a line.
Choose a real row vector $w_0$ nonzero on that line. The map
\[
 T_w:x\longmapsto (\ell(x),wx,\pi(x))\in\mathbb R^2\oplus F
\]
is invertible at $w=w_0$. There are an open neighborhood $U$ of $w_0$ and
$K<\infty$ such that $\|T_w^{-1}\|\le K$ for every $w\in U$.

The projected Parikh vectors of factors of $W$ are uniformly bounded.
Indeed, successively removing whole substitution blocks from a prefix $z$
gives a finite expansion
\[
 \Psi(z)=\sum_{j\ge0}A^j\Psi(p_j),
\]
where each $p_j$ is a proper prefix of an image of the powered substitution.
The parent prefix has at most half the previous length, so this expansion
terminates. Only finitely many $p_j$ are possible, $\pi$ commutes with $A$,
and $\sum_j\|A^j|_F\|<\infty$. Hence projected prefixes are bounded. A factor is
a difference of two prefixes, so there is $B$ such that
$\|\pi\Psi(u)\|\le B$ for every factor $u$.

For two factors put $d=\Psi(u)-\Psi(v)$. Then $\|\pi d\|\le2B$.
If $\ell(d)=wd=0$ for $w\in U$, the inverse bound gives $\|d\|\le2KB$.
There are finitely many nonzero integer vectors $d$ in this ball with
$\ell(d)=0$. Choose a rational $w\in U$ outside their finitely many
hyperplanes $wd=0$, and outside the hyperplanes $w_i=w_j$ for $i\ne j$.
Multiplying by a common denominator and adding a sufficiently large integer
multiple of $\ell$ makes the weights distinct positive integers. These two
operations preserve vanishing on vectors with $\ell(d)=0$.

Consequently, factors with equal length and equal weighted sum have equal
Parikh vectors. An additive square in the weighted word would therefore be an
abelian square in $W$, a contradiction. Finally let $f(0)=0$ and let
$f(n+1)-f(n)$ be the $n$th weight. Its bounded increments make it Lipschitz,
and $f(n)+f(n+2r)=2f(n+r)$, $r>0$, is precisely the excluded additive square.
\end{proof}

\paragraph{Exact spectral filtering.}
The search uses integer characteristic polynomials and a Cayley transform,
rather than approximate eigenvalues. If $q$ has degree $d$, set
\[
 R(s)=(1-s)^d q\!\left(\frac{1+s}{1-s}\right).
\]
Roots outside the unit disk correspond to roots of $R$ in the open right
half-plane. Roots $1$ and $-1$ of $q$ are detected directly. In a nonsingular
Routh table the signs of its first column are those of
\[
 a_0,\quad \Delta_1,\quad
 \Delta_2/\Delta_1,\ldots,\Delta_d/\Delta_{d-1},
\]
where $a_0$ is the leading coefficient of $R$ and $\Delta_j$ its leading
Hurwitz determinants. Their sign changes count right-half-plane roots.
All determinants in the program are computed with integers.

Singular tables are handled by exact rational shifts as follows.
If $R$ is an integer polynomial of degree $d$, height $H$, and nonzero
leading and constant coefficients, put
\[
 K=(2d)!\,(H2^{d+1})^{2d},\qquad J=d^2(d+1)+1.
\]
The nonzero polynomial
$T(t)=\operatorname{Res}_s(R(s),R(t-s))$ has integer coefficients of absolute
value at most $K$. To see the height bound, each entry of its Sylvester matrix
has coefficient $\ell^1$ norm at most $H2^{d+1}$; expand the determinant and
use the $(2d)!$ possible permutation terms. Its roots include
$\alpha+\overline\alpha=2\operatorname{Re}\alpha$ for every root $\alpha$ of $R$.
After removing any power of $t$, Cauchy's reciprocal root bound shows that
every nonzero root of $T$ has absolute value at least $1/(1+K)$. Thus
\[
 \operatorname{Re}\alpha\ne0
 \quad\Longrightarrow\quad
 |\operatorname{Re}\alpha|\ge\frac1{2(1+K)}.
\]
Choose an integer $Q>4J(1+K)$. For some $1\le j\le J$, both Routh tables
of $R(s+j/Q)$ and $R(s-j/Q)$ are nonsingular. Indeed each size-$k$ Hurwitz
determinant is a nonzero polynomial of degree at most $kd$ in the shift:
nontriviality follows by moving all roots into one open half-plane.
Across both signs these polynomials have total degree at most $d^2(d+1)$.
The two right-half-plane counts respectively exclude and include the
imaginary-axis roots of $R$; their difference counts those roots with
multiplicity. The chosen shifts cannot cross any other root by the bound above.

\paragraph{Completed finite search.}
The starting substitution, due to Rao and Rosenfeld, is
\[
 a\mapsto ace,\quad b\mapsto adf,\quad c\mapsto bdf,\quad
 d\mapsto bdc,\quad e\mapsto afe,\quad f\mapsto bce.
\]
Its known abelian-square-free fixed point is a positive control, and its
incidence characteristic polynomial is $x^3(x-3)(x^2-3)$.
The program changes at most five of its eighteen output symbols, keeping
the initial $a$ in $h(a)$ fixed. It therefore enumerates exactly
\[
 \sum_{j=0}^{5}\binom{17}{j}5^j=20\,913\,486
\]
substitutions. Of these, $5729$ pass the abelian-square check through
$h^4(a)$, of length $81$; $5715$ of those have primitive incidence matrices.
Exactly $73$ of those meet the two-expanding-direction condition and have no
unit-modulus eigenvalues. Every one has an abelian square by position $203$,
with positions numbered starting at one. Thus no substitution in this finite
family satisfies all the theorem's hypotheses. This says nothing about more
distant substitutions, other image lengths, other alphabet sizes, or the
existence question itself.

The suffix checker was compared against direct Parikh counts on every
length-six word over three letters. The unmodified six-letter substitution
passed a $6561$-letter positive-control check. The Routh counter was checked
on products of known linear and quadratic factors, including repeated and
reciprocal roots. The projection proof and search logic received separate
fresh-context reviews; the search reviewer did not independently rerun the
complete radius-five enumeration.

\paragraph{Files and remaining gap.}
The executable is \texttt{green15-morphism-search.mjs}; run
\texttt{node green15-morphism-search.mjs 5}.
The completed output is \texttt{green15-morphism-results-5-routh-exact.json}.
The main unresolved task is to construct an infinite square-free word with
the requisite spectral structure, or to find a different argument settling
Green's Problem 15. No such construction or impossibility proof is supplied.

\paragraph{Complete short-image exclusion.}
A second completed computation covers all primitive substitutions on eight
letters with four images of length one and four of length two, up to renaming,
subject to the spectral hypothesis of the theorem. A cycle consisting only of
length-one letters would contradict primitivity. Their maps therefore form
forests rooted at the four other letters. Enumerating all $8^4$ maps gives
$2048$ acyclic maps, in $19$ classes under independent permutations of the
four letters of each type. Equivalently, distributing four nonroot vertices
among four unlabeled rooted components gives $9+4+3+2+1=19$ shapes.

For each representative there are $64^4$ assignments of the two-letter
images. The program assigns an image only when an iterate of the starting
letter requests it. A square in a determined word prefix rejects every
completion. A closed reachable subalphabet omitting letters rejects every
completion by failure of primitivity. Otherwise all $64$ assignments of the
requested image are explored. Each pruned node accounts for $64^u$ completions,
where $u$ is the number of unassigned images. Thus the checked coverage is
\[
 19\cdot64^4=318\,767\,104.
\]
The computation inspected $33\,285\,888$ complete assignments after pruning;
$1948$ candidates reached the longer test after the initial factor, primitivity,
and spectral filters. None survived. The largest reported valid prefix of a
tested generated factor was $207$, followed by a square ending at position $208$.
The known eight-letter Rao--Rosenfeld example, canonically renamed, reached its
expected leaf exactly once as a control.

The iterates in this search need not be nested prefixes. For primitive $h$,
every $h^k(b)$ is a factor of a fixed point $W$ of a prolongable power $h^t$:
choose $m$ large enough that $h^{tm-k}(a)$ contains $b$, and use the occurrence
of $h^k(b)$ in the prefix $h^{tm}(a)$ of $W$. This justifies the square pruning.
The complete record is \texttt{green15-forest-results.json}, obtained with
\texttt{node green15-morphism-search.mjs 0 --forests}.

\begin{lemma}[Modular normal form]
Fix $t\ge1$, $L=4^t$, and $M=3^t$. There exists an infinite additive-square-free
word over a finite integer alphabet if and only if one exists whose prefix sums
$S$ satisfy
\[
 S(n)\equiv F_t(n\bmod L)\pmod M,
 \qquad
 F_t(n)=\sum_{j=0}^{t-1}3^j\mathbf1_{\{n_j\ge2\}},
\]
where $n_j\in\{0,1,2,3\}$ are the base-four digits of $n\in[0,L)$.
\end{lemma}
\begin{proof}
The cyclic pattern $F_1=(0,0,1,1)$ has no three-term arithmetic progression
in its graph modulo $(4,3)$ with nonzero first-coordinate step; checking the
four starting positions and steps $1,2,3$ proves this. Inductively, if
\[
 F_t(n)+F_t(n+2d)-2F_t(n+d)\equiv0\pmod {3^t},
\]
with arguments reduced modulo $4^t$, reduction modulo $3$ forces $4\mid d$.
Writing $n=4q+r$ then divides the congruence by $3$ and gives the corresponding
congruence for $F_{t-1}$ at $q,q+d/4,q+2d/4$. Hence $4^t\mid d$.

Put $F_t(L)=F_t(0)=0$ and $c_r=F_t(r+1)-F_t(r)$. Code an input integer $a$ by
\[
 c_0,c_1,\ldots,c_{L-2},c_{L-1}+Ma.
\]
This sends a finite input alphabet to a finite output alphabet. The output
prefix sums are
\[
 S_{\rm out}(Ln+r)=M S_{\rm in}(n)+F_t(r),\qquad 0\le r<L.
\]
An output additive square must have half-length divisible by $L$, by the
cyclic property just proved. All three cut positions then have the same
offset $r$; the displayed formula gives an input additive square, a
contradiction. The converse is immediate by forgetting the marker condition.
\end{proof}

This is a normalization of the full word-existence question. The searches
impose the additional hypothesis of a finite substitution; the lemma does
not establish that hypothesis. Nor does iterating this particular code keep
its alphabet finite: its exceptional output letter is
$a\mapsto Ma-(M-1)/2$. A finite orbit would force its fixed point $a=1/2$,
which is not an integer.

\paragraph{Affine synthesis and its current limits.}
The marker searches use $L$-uniform substitutions with
\[
 h(a)_r=c_r+Mz_{a,r}\in[-K,K],\qquad
 \sum_r z_{a,r}=a,\qquad h(0)_0=0.
\]
Thus image sums equal $Ma$. The prefix is determined by
$W(0)=0$ and $W(n)=h(W(\lfloor n/L\rfloor))_{n\bmod L}$ for $n>0$.
All additive-square half-lengths must be multiples of $L$.

The local Z3 solver uses dependency-conditioned linear cuts. A concrete
square determines the image-entry variables used at its positions. Fixing
the image-entry variables needed to determine their parent letters fixes
which entries are selected. Under that context, equality of block sums is
one integer linear equation in the selected $z_{a,r}$. The solver learns
the disjunction that the parent context changes or this equation is nonzero.
The argument and its implementation received a fresh-context review with
no substantive issue reported.

For $t=1,K=3$, the solver returned unsatisfiable after $26$ learned cuts,
reproducing an independent exhaustive backtracking exclusion whose longest
valid prefix was $42$. No independent solver proof certificate is supplied.
For $t=2,K=18$, a one-cut-per-model run checked $404$ models, reached a
$422$-letter valid prefix, and returned \texttt{unknown} on its next solver
call. That outcome leaves the family unresolved. Separate bounded parabolic
image searches reached a $1954$-letter prefix and hit their node limits;
they also provide no infinite construction or exhaustive exclusion.

The relevant executables are \texttt{green15-marker-z3.mjs},
\texttt{green15-marker-synthesis.mjs}, and \texttt{green15-parabolic-search.mjs}.
Z3's official JavaScript package, version 5.3.3, is installed locally under
\texttt{.research-tools}; all generated candidates are checked with exact
integer sums outside the solver. \textbf{The full problem remains unsolved.}

\begin{thebibliography}{9}
\bibitem{rr} M. Rao and M. Rosenfeld,
\emph{Avoiding two consecutive blocks of same size and same sum over
$\mathbb Z^2$}, \url{https://arxiv.org/abs/1511.05875}.
\bibitem{green} B. Green, \emph{100 open problems}, Problem 15,
\url{https://people.maths.ox.ac.uk/greenbj/papers/open-problems.pdf}.
\bibitem{routh} X. Chen, \emph{Stability Analysis: the Routh--Hurwitz Approach},
\url{https://faculty.washington.edu/chx/teaching/fbintro/lec10-routh-hurwitz-stability/}.
\end{thebibliography}
\end{document}
